\section{Summary and Conclusions} \label{sec:conclusion}\noindent
This study introduced a simplified neural metamodel specifically designed to address the challenges of spatiotemporal, multi-field responses in geotechnical systems. By decoupling spatial and temporal learning while maintaining a unified latent representation, the proposed architecture links low-dimensional static soil parameters to complex, evolving outputs, including displacement fields, stress distributions and deformation vectors, within a single multi-output surrogate.

Through two benchmark problems, an offshore strip footing and a triaxial loading test, the metamodel demonstrated good predictive accuracy, robustness and efficiency. In the footing case, it captured both vector- and field-type responses across multiple loading stages with a test-set global normalised absolute error (GNAE) of about 2.3\%, while in the triaxial test the GNAE was about 1.3\%; For both examples, MAE and relative $L_2$ errors showed consistent trends for the investigated examples. Importantly, the metamodel was not only assessed in forward-prediction mode but also embedded in a sequential Bayesian calibration framework with multiple outputs, where it enabled efficient likelihood evaluations and led to a substantial reduction in parameter uncertainty as monitoring information accumulated.

The simplified architecture thus provides a practical balance between accuracy, computational efficiency and engineering applicability. Its modular, multi-output design facilitates extension to additional response quantities, heterogeneous soil conditions and more complex loading scenarios, and its compatibility with existing probabilistic tools makes it a promising surrogate for real-time prediction and uncertainty quantification in geotechnical practice. Future work may focus on enhancing generalisation to a wider range of constitutive models, incorporating spatial variability in ground conditions, mesh/geometry sensitivity and scaling the approach to large-scale infrastructure systems.